\section{Low-F multistandarding}
\label{app:low-f}

Main theoremは意図的にhigh-cost regionとintermediate-cost regionを比較する。完全性のため、本Appendixでは対称モデルにおける検証済みの十分なlow-$F$ regionを記録する。この結果は選択的侵食メカニズムではない。すべての企業がすべての関連標準をサポートすると、正式レジームの違いは重複した固定採用費用だけになる。

\begin{equation}
0<F<\min\{T_U,T_A\}.
\label{eq:appendix-low-f-condition}
\end{equation}
とする。$T_W>T_U$ なので、condition \eqref{eq:appendix-low-f-condition} は $F<T_W$ も含意する。SUの下では、$F<T_A<2T_A$ なので域外国企業は厳格にbloc standardを採用する。2つの加盟企業についても、他方が採用していないときは $F<T_U$、採用しているときは $F<T_A$ なので、どちらの場合も域外国標準の採用が厳格に有利である。したがって3社すべてが2つの正式標準をサポートする。

SWの下では、すべてのtarget standardについて、2つの外国企業はライバルが採用しているかどうかにかかわらず採用を厳格に選好する。$F<T_W$ かつ $F<T_A$ だからである。したがって各企業はnative standardに加えて両方のforeign standardをサポートする。国際標準化はすでに単一の共通標準を持ち、私的採用を必要としない。

Universal multistandardingにより、すべてのdestinationがcomplete compatibility product marketとなる。したがって消費者余剰とgross operating profitはISと同じである。残る唯一の違いは、各国内企業が支払うadditional-standard fixed costの数である。SUでは各企業が1つの追加標準をサポートするのに対し、SWでは2つをサポートする。したがって、すべての国について
\begin{equation}
\boxed{
W_i^{SU}=W_i^{IS}-F,
\qquad
W_i^{SW}=W_i^{IS}-2F.
}
\label{eq:appendix-low-f-welfare}
\end{equation}
である。$F>0$ なので、この十分なlow-cost regionでは、すべての国が非IS formal regimeのどちらの下でもISより低い継続利得を得る。これはtheorem ledgerに記録されたfrozen low-$F$ formal-regime comparisonである。この結果を用いてTheorem~\ref{thm:formal-coalition-stability} を記載されたhigh-およびintermediate-$F$ regionの外へ拡張することはしない。

本Appendixの結果を、$F$ に関するglobal monotonicity statementへ外挿してはならない。その経済学はheadline intermediate-cost resultとは異なる。intermediate regionでは、域外国企業のみの採用がregional exclusion rentを選択的に取り除く一方、reciprocal incompatibilityが残る。十分なlow-$F$ regionでは、effective compatibilityはすでにuniversalであり、残るformal-regime differenceは重複したprivate adoption costの会計だけである。したがってMain contributionは、Sections~\ref{sec:selective-erosion} および \ref{sec:coalition-stability} で確立したselective-erosion transitionのままである。
